\documentclass[10pt,a4paper]{amsart}
\usepackage[margin=19mm]{geometry}
\usepackage{amsmath,amssymb,mathtools}
\usepackage[scr=rsfs,cal=euler]{mathalfa}
\usepackage{xfrac}
\usepackage[unicode,hidelinks,bookmarks=false]{hyperref}
\newcommand{\ie}{i.e.\ }
\newcommand{\cf}{cf.\ }
\newcommand{\resp}{respectively\ }
\newcommand{\Subsection}{\subsection}
\providecommand{\nobreakdash}{\nobreak\hbox{-}}
\setlength{\emergencystretch}{2em}
\multlinegap0pt
\usepackage{amssymb}
\usepackage{mathtools}
\usepackage[normalem]{ulem}
\newtheorem{lemma}{Lemma}
\newtheorem{setting}{Setting}
\newtheorem{theorem}{Theorem}
\newtheorem{prop}{Proposition}
\newcommand{\bC}{{\mathbb C}}
\title{Bergman spaces on algebraic curves}
\author{L\'aszl\'o Koltai, Alexander A. Kubasch, R\'obert Sz\H{o}ke}
\date{Source: arXiv:2504.02341v2 (CC BY 4.0)}
\begin{document}
\maketitle

\noindent\emph{Source and licence.} Bergman spaces on algebraic curves. Version arXiv:2504.02341v2 (CC BY 4.0), distributed by arXiv under the Creative Commons Attribution 4.0 International licence, http://creativecommons.org/licenses/by/4.0/, as recorded in the arXiv licence metadata for this exact version and shown on the abstract page https://arxiv.org/abs/2504.02341v2. Full licence notice: \texttt{licenses/arxiv-2504.02341-CC-BY-4.0.txt}.\par\smallskip

\noindent\textit{Editorial scope.} Counted unit: the article's theorem thm:Phiisom (source0.tex lines 370--392). The article's complete original proof of it is delivered with it (source0.tex lines 414--486, one \texttt{proof} environment of 73 lines, which the article places away from the statement). No mathematics is written, repaired or re-derived: the statement and the proof are the article's own lines, in the article's own notation, and no retained line is substituted or re-flowed.  Retained uncounted as the source-local context the proof needs: lines 210--222; lines 356--368.  Nothing was omitted from any retained span, so the package carries no omission record.  No retained line contains a citation, so the wrapper carries no bibliography and no bibliography entry.  No retained line contains a figure environment, an image-inclusion command or drawing code, so no image is delivered and the package declares no asset dependency.  Editorial formatting only, in addition: an AMS \texttt{amsart} preamble that restates the article's own declarations and macros used by the retained lines, namely the environments \texttt{lemma}, \texttt{setting}, \texttt{theorem}, \texttt{prop} on separate counters, together with the standard packages those lines need.  The retained environments print in this document as Lemma 1, Setting 1, Theorem 1, Proposition 1 in order of appearance, whereas the article prints the same items with the numbering of its own full document.  The complete native source of the article as distributed in the arXiv e-print for this exact version is bundled unchanged as \texttt{sources/175884/source0.tex}.
\begin{lemma}\label{lem:korlap}
     Let $0 \in U\subset\bC$ be open and  $m\in \mathbb{Z}$.
     \begin{enumerate}
        \item Suppose $m < 0$. Then the map
         $$\Phi_- : A^2(U,\mu_m) \to A^2(U), \quad f(z) \mapsto z^mf(z)$$
         is a surjective linear isometry. In this case we also have that
            $$ A^2(U,\mu_m)=A^2(U^*,\mu_m).$$
         \item Suppose $m > 0$. Then the map
         $$\Phi_+ : A^2(U,\mu_m) \to A^2_m(U,0), \quad f(z) \mapsto z^mf(z)$$
         is a surjective linear isometry.
     \end{enumerate}
     
\end{lemma}

\begin{setting}\label{S:input}
 Let $M$ be a Riemann surface, $\omega$ a positive continuous volume form on $M$,
 $(E,h)\rightarrow M$ an Hermitian holomorphic vector bundle with $h$ continuous and
 $D$ a divisor on $M$. Let
$D_+:=\sum_{m_p>0}m_pp$ and  
  $D_-:=\sum_{m_p<0}m_pp$.
    Then $D=D_++D_-$. Denote by  $L_{D_\pm}$ and  $L_D$ the associated holomorphic line bundles.  $L_D$
    is isomorphic to $L_{D_+}\otimes L_{D_-}$. Let $s_+:M\rightarrow L_{D_+}$ (resp. $s_-:M\rightarrow L_{D_-}$) be a holomorphic (resp. meromorphic) section with $\text{div}(s_\pm)=D_\pm$.
    Let $\kappa_+$ (resp. $\kappa_-$) be a continuous Hermitian  metric on $L_{D_+}$ (resp. on $L_{D_-}$).
    Then $\kappa=\kappa_+\otimes \kappa_-$ is a continuous Hermitian metric on $L_{D_+}\otimes L_{D_-}$.
     Let $K_\pm=\kappa_\pm(s_\pm,s_\pm)$ and
     $\rho_D:=\kappa(s_D,s_D)\omega$  a $D$-volume form.  
\end{setting}

\begin{theorem}\label{thm:Phiisom}
Using the definitions and notations of Setting \ref{S:input} the following hold:
 
 \begin{enumerate}
     \item 
  Suppose  $D\le0$. Then $s\in A^2(M,\rho_D,E,h)$ implies that $ 
 -D\le \text{\normalfont div}(s).$ Furthermore, the map
 $$
 \Phi_-:A^2(M,\rho_D,E,h)
 \rightarrow A^2(M,\omega,E\otimes L_D, h\otimes\kappa),\quad s\mapsto s\otimes s_D
 $$
 is a surjective linear isometry.

\item Suppose  $0\le D$. The 
 map
$$\Phi_+:A^2(M,\rho_D,E,h)
 \rightarrow A^2_D(M,\omega,E\otimes L_D, h\otimes\kappa),\quad s\mapsto s\otimes s_D
 $$
 is a surjective linear isometry.
 \item Finally, for a general divisor $D \in \text{\normalfont Div}(M)$, the map $$\Phi : A^2(M,\rho_D,E,h) \to A^2_{D_+}(M \setminus \text{\normalfont supp}(D_-), K_-\omega, E\otimes L_{D_+}, h \otimes \kappa_+ ), \quad s \mapsto s\otimes s_{D_+}\big|_{M \setminus \text{\normalfont supp}(D_-)}$$ is a surjective linear isometry.
 
 \end{enumerate}
\end{theorem}

\begin{proof}[Proof of Theorem \ref{thm:Phiisom}] First we show that both maps $\Phi_{\pm}$ are linear isometries and that they map into the appropriate Bergman space.
 Let $s$ be a holomorphic section of $E$. Then  
\begin{equation}\label{eq:isometry}
\begin{split}
\|s\|^2_{L^2(\rho_D)}=\int_Mh(s,s)\rho_D=
\int_M h(s,s)\kappa(s_D,s_D)\omega=\\  =\int_M h\otimes \kappa(s\otimes s_D,s\otimes s_D)\omega=
\|s\otimes s_D\|^2_{L^2(\omega)}.\hspace{10mm}
\end{split}   
\end{equation} 
Hence
\begin{equation}\label{eq:equiv}
\|s\|^2_{L^2(\rho_D)}<\infty\Longleftrightarrow
\|s\otimes s_D\|^2_{L^2(\omega)}<\infty
\end{equation}
Suppose now that $D\le 0$.
 Since isolated $L^2$ singularities of holomorphic functions are removable, \eqref{eq:equiv} yields  that if $s\in  A^2(M,\rho_D,E,h)$, the section $s\otimes s_D$ extends holomorphically to the points of $\text{supp}(D)$ and so
 \begin{equation}\label{E:D<0}
 0\le \text{div}(\Phi_-(s))= \text{div}(s\otimes s_D)=\text{div}(s)+D\Longrightarrow
 -D\le \text{div}(s).
\end{equation}
 Equation~\eqref{eq:isometry} implies that the map $\Phi_-$   is a  linear isometry into the Bergman space $A^2(M,\omega,E\otimes L_D, h\otimes\kappa)$. 
 
 
 If $0\le D$, the section $s_D$ is holomorphic
 and $s\otimes s_D$ is as well.
 Furthermore 
 $$\text{div}(\Phi_+(s))=\text{div}(s\otimes s_D)=\text{div} (s)+D\ge D,$$
 since $s$ is holomorphic. This together with
 equation~\eqref{eq:isometry}
yields that $\Phi_+$ indeed maps into  $A^2_D(M,\omega,L\otimes L_D, h\otimes\kappa)$ and that it is a linear isometry.
What is left  to show is that both maps $\Phi_-$ and $\Phi_+$ are surjective. 

Let  $S\in H^0(M,E\otimes L_D)$. Since the bundle $L_D$ is trivial over $M\setminus{\text{supp}(D)}$, there exists a holomorphic section $s$ of $\left.E\right|_{M\setminus{\text{supp}(D)}}$ with
$s\otimes\left.s_D\right|_{M\setminus{\text{supp}(D)}}=\left.S\right|_{M\setminus{\text{supp}(D)}}$.
\begin{prop} \label{prop:smerom}
    $s$ is a meromorphic section of $E\rightarrow M$.
\end{prop}
\begin{proof}
 Let $p\in \text{supp}(D)$. Choose a small open neighborhood $U_p$ of $p$ and let $e_1,\dots,e_r$ and $e$ be holomorphic frames of $E$ and $L_D$ respectively. Then $\left.S\right|_{U_p}=\sum g_je_j\otimes e$ for some $g_j\in\mathcal O(U_p)$ and $\left.s_D\right|_{U_p}=he$ where $h$ is a meromorphic function on $U_p$. On $U^*_p:=U_p\setminus\{p\}$, $s=\sum f_je_j$ for some $f_j\in\mathcal O(U^*_p)$.
 Therefore
 $$
\sum_j g_j\left.(e_j\otimes e)\right|_{U^*_p}= \Big(\sum_j f_je_j\Big)\otimes\left.(he)\right|_{U^*_p}=
\sum_j (hf_j)e_j\otimes \left.e\right|_{U^*_p}
 $$
 and hence $f_j=g_j/h$ is meromorphic
 in $U_p$ for all $j=1,\dots, r$.
\end{proof}
\noindent It follows from Proposition \ref{prop:smerom} that
\begin{equation}\label{eq:div}
 0\le\text{div}(S)=\text{div}(s)+\text{div}(s_D)=\text{div}(s)+D.   
\end{equation}
Now, let $D<0$ and $S\in A^2(M,\omega, E\otimes L_D,h\otimes\kappa)$. Then 
\eqref{eq:div} shows that $s$ is holomorphic and \eqref{eq:isometry}  yields that $s\in A^2(M,\rho_D,E,h)$ 
 proving the surjectivity of $\Phi_-$.
 


Now, suppose that $D\geq0$ and 
$S\in A^2_D(M,\omega,L\otimes L_D, h\otimes\kappa)$.
Then \eqref{eq:div} shows again that $s$ is a holomorphic section and \eqref{eq:isometry} that
 $s\in A^2(M,\rho_D,E,h)$ 
 proving the surjectivity of $\Phi_+$.

 Finally, let $D \in \text{Div}(M)$ be an arbitrary divisor.  Isolated $L^2$ singularities of holomorphic functions are removable, hence the restriction map
 $$
 \iota: H^0(M,E)\rightarrow H^0(M\setminus\text{supp} (D_-),E),\quad s\mapsto 
 \left.s\right|_{M\setminus\text{supp}(D_-)}
 $$
 induces an onto isometry
 $$
\iota^*:A^2(M,\rho,E,h)=A^2(M,K_+K_-\omega,E,h)\rightarrow A^2(M\setminus\text{supp}(D_-), K_+K_-\omega, E,h)$$
by Lemma \ref{lem:korlap}. Hence, applying part (2) to the effective divisor $D_+$ shows that $\Phi := \Phi_+ \circ \iota^*$ is a surjective linear isometry as well.
\end{proof}

\end{document}
